%% ENGR489 Final Report -- high-level outline for the supervisors: each section lists what it covers.
%% Preamble and front matter follow engr489_report_template.tex (IEEEtran journal).
%%
%% Handbook (Ch. 5): max 14 pages inclusive of references, figures, and tables; do not change
%% font settings; 2-minute video submitted separately; appendix (if any) is a separate,
%% unmarked document. Page budgets per section below assume 14 pages in total.
%%
%% Style guide reminders (Honours_Report_Style_Guide_Engineering.pdf): title-case headings,
%% NZ English, no first person, define acronyms on first use, one to nine in words / 10+ as
%% numerals, captions end with a full stop (figures below, tables above), every figure and
%% table referenced in text, IEEE citations.
%%
%% Detailed figures (test counts, coverage, latency, SUS scores, graph statistics) are in
%% documentation/backend-fixes.md, documentation/spec.md and this file's git history.

\documentclass[lettersize,journal]{IEEEtran}
\usepackage{amsmath,amsfonts}
\usepackage{array}
\usepackage[caption=false,font=normalsize,labelfont=sf,textfont=sf]{subfig}
\usepackage{textcomp}
\usepackage{stfloats}
\usepackage{url}
\usepackage{verbatim}
\usepackage{graphicx}
\usepackage{cite}
\hyphenation{op-tical net-works semi-conduc-tor IEEE-Xplore}

\begin{document}

\title{Hunting Mr.~X: Wellington Edition}

\author{Ming Bao
\thanks{This project was supervised by Jens Dietrich (primary) and Stuart Marshall.}
}

\markboth{ENGR 489 (ENGINEERING PROJECT) 2026}%
{Bao: Hunting Mr.~X: Wellington Edition}

\maketitle

%%%%%%%%%%%%%%%%%%%%%%%%%%%%%%%%%%%%%%%%%%%%%%%%%%%%%%%
\begin{abstract}
\begin{itemize}
	\item Problem: hidden-movement board games need everyone in the same room, and existing digital versions use fictional boards.
	\item Solution: a browser-based, real-time multiplayer game played on a map of Wellington's real transport network.
	\item Evaluation: automated testing, a map-library benchmark, and two rounds of user testing with a usability questionnaire.
	\item Outcome: a deployed, playable system suitable for the university Open Day.
\end{itemize}
\end{abstract}

\begin{IEEEkeywords}
Multiplayer game, hidden movement, real-time web application, WebSocket, map overlay, usability evaluation.
\end{IEEEkeywords}

%%%%%%%%%%%%%%%%%%%%%%%%%%%%%%%%%%%%%%%%%%%%%%%%%%%%%%%
\section{Introduction} % ~1.5 pages
\begin{itemize}
	\item Motivation: why hidden-movement games suit online play, and the gap this project fills (no browser version on a real city map).
	\item Secondary goal: a showpiece for the university Open Day.
	\item Aims and measurable goals, which the Evaluation section refers back to.
	\item Contributions, and a note on naming: the mechanics are used with the publisher's permission, and the original brand name is not used.
	\item Key findings and a roadmap of the report.
\end{itemize}

%%%%%%%%%%%%%%%%%%%%%%%%%%%%%%%%%%%%%%%%%%%%%%%%%%%%%%%
\section{Related Work} % ~1.5 pages
\begin{itemize}
	\item The original board game and existing digital versions of it.
	\item Location-based games that use real-world maps.
	\item Keeping hidden information secret in online multiplayer games.
	\item Map libraries, and live routing services versus pre-computed routes \cite{placeholder_web}.
	\item Real-time communication for web applications.
	\item Measuring usability with the System Usability Scale (SUS).
	\item How this project addresses the gaps these leave.
\end{itemize}
% Citations still to find: the Ravensburger rulebook; specific digital versions (the progress-presentation
% table is a starting point); a game-networking source; map-library docs and licences; the OpenStreetMap
% licence; Brooke (SUS) and Sauro (the 68 benchmark).

%%%%%%%%%%%%%%%%%%%%%%%%%%%%%%%%%%%%%%%%%%%%%%%%%%%%%%%
\section{Requirements and Constraints} % ~1 page
\begin{itemize}
	\item Functional requirements: the game rules the system has to enforce.
	\item Non-functional requirements: responsiveness, several games at once, and browser support.
	\item Constraints: no paid per-request services, no database, a solo project on a small budget, and self-hosting.
	\item A table linking each requirement to where it is evaluated.
\end{itemize}

%%%%%%%%%%%%%%%%%%%%%%%%%%%%%%%%%%%%%%%%%%%%%%%%%%%%%%%
\section{Design} % ~3 pages
\subsection{Architecture}
\begin{itemize}
	\item A client-server design in which the server is the single source of truth.
	\item The alternatives considered and why they were not chosen.
\end{itemize}
\subsection{Game State Machine and Turn Flow}
\begin{itemize}
	\item How a game moves from the lobby through play to its end.
	\item What happens when a player leaves or stops responding.
\end{itemize}
\subsection{Domain Model}
\begin{itemize}
	\item How players, tickets and rules are modelled, with the rules kept apart from the web framework so they can be tested directly.
\end{itemize}
\subsection{Hidden Information and Role-Filtered State}
\begin{itemize}
	\item How each player is sent only what they are allowed to see, so Mr~X's position stays secret.
\end{itemize}
\subsection{Map Data Design}
\begin{itemize}
	\item Why the Wellington map is a pre-computed graph rather than live routing.
	\item How the map library was chosen.
\end{itemize}
\subsection{User Interface Design}
\begin{itemize}
	\item The main screens, and the changes made in response to user feedback.
\end{itemize}
% Add mockups or the Figma reference here, if used.
\subsection{Real-World Issues}
\begin{itemize}
	\item Ethics, sustainability, safety, intellectual property and privacy.
\end{itemize}
% Check the ethics wording (voluntary student participants, course coordinator approval, vouchers).

%%%%%%%%%%%%%%%%%%%%%%%%%%%%%%%%%%%%%%%%%%%%%%%%%%%%%%%
\section{Implementation} % ~3 pages
\begin{itemize}
	\item Technology choices and the reasons for them.
	\item How the server enforces the rules and runs several games at once.
	\item How the browser client shows the map, the game state and live updates.
	\item Building the Wellington map, the tooling made for it, and the problems met along the way.
	\item Deployment, and keeping the API documentation in step with the code.
	\item The development process, including how AI-assisted development was used and checked.
	\item Screenshots of the main screens.
\end{itemize}
% Check how the course expects AI use to be disclosed.

%%%%%%%%%%%%%%%%%%%%%%%%%%%%%%%%%%%%%%%%%%%%%%%%%%%%%%%
\section{Evaluation} % ~3.5 pages
\subsection{Map Library Benchmark}
\begin{itemize}
	\item How three map libraries were compared, the results, and the limits of the benchmark.
\end{itemize}
\subsection{Functional Correctness}
\begin{itemize}
	\item The layers of automated tests and what each one checks.
	\item Regression tests for the bugs found in the backend review.
	\item Coverage, mutation and property-based testing as checks on the tests themselves.
\end{itemize}
\subsection{Performance and Multiplayer Synchronisation}
\begin{itemize}
	\item Update latency and map rendering, measured against the targets from the Introduction.
\end{itemize}
% To do: measure map render time on the real map.
\subsection{User Evaluation}
\begin{itemize}
	\item Two rounds of user testing and the SUS questionnaire.
	\item The main themes in the feedback and the changes they led to.
	\item Limitations of the study.
\end{itemize}
% Possibly a before/after comparison if a third round is run. Decide how to read the side-balance
% question (check the form's scale).
\subsection{Requirements Traceability}
\begin{itemize}
	\item Each goal from the Introduction marked as met, partly met or not met.
\end{itemize}

%%%%%%%%%%%%%%%%%%%%%%%%%%%%%%%%%%%%%%%%%%%%%%%%%%%%%%%
\section{Discussion} % ~0.5 page
\begin{itemize}
	\item What worked well and what was hard.
	\item Lessons learned, including a security weakness found in review and how it was fixed.
	\item Trade-offs revisited, and threats to validity.
\end{itemize}

%%%%%%%%%%%%%%%%%%%%%%%%%%%%%%%%%%%%%%%%%%%%%%%%%%%%%%%
\section{Conclusions and Future Work} % ~0.5 page
\begin{itemize}
	\item What was delivered and the main results.
	\item Future work for a following student: a tutorial, configurable lobbies, saved games and accounts, mobile support, automated map generation, balance testing with bot players, and a larger user study.
\end{itemize}

%%%%%%%%%%%%%%%%%%%%%%%%%%%%%%%%%%%%%%%%%%%%%%%%%%%%%%%
\section*{Acknowledgments}
\begin{itemize}
	\item Supervisors, the course coordinator and the user-testing participants.
\end{itemize}
% Add the friend or teammate who helped with the map, if any.

%%%%%%%%%%%%%%%%%%%%%%%%%%%%%%%%%%%%%%%%%%%%%%%%%%%%%%%
%% Appendix (repository link, openapi.yaml, full spec, diagrams, anonymised SUS responses --
%% drop the email column -- map-creator notes, supervisor meeting log) goes in a SEPARATE
%% document per the handbook; it is not marked and does not count toward the 14 pages.

% Replace with a BibTeX .bbl or \bibitem entries in IEEE style. Candidates to cite:
%  - Ravensburger rules (original board game)
%  - MapLibre GL, Leaflet, Google Maps documentation
%  - Spring Framework WebSocket/STOMP documentation; Vue, Vite, Tailwind
%  - OpenStreetMap / Metlink data sources for the graph (confirm what was used)
%  - J. Brooke, ``SUS: a quick and dirty usability scale''; J. Sauro, MeasuringU SUS benchmark
%  - Game-networking and hidden-information references (to find)
\bibliographystyle{IEEEtran}
\bibliography{references}

\end{document}
